


\section{Predictors}
\label{sec:rh_predictors}
\begin{itemize}
\item \textbf{Trailing averages on the full grid.} Following \citet{Corsi2009}, every predictor is a trailing average
  of a transformed series $\tilde u$ over $k$ samples, lagged by one sample:
  \[
    \bar u^{(k)}_{t,i} = \frac{1}{k} \sum_{j=1}^{k} \tilde u_{t,i-j}, \qquad
    k \in \mathcal{K} = \{1, 5, 25, 125, 625, 3125\},
  \]
  \rhdmNHorizonsPooled{} horizons in base 5. A sample is one interval of the full grid $\mathcal{G}$, overnight
  sessions included: the models are estimated on the \rhorNBars{} regular-hours bars of each day
  (Section~\ref{sec:rh_methods_setup}), but the averages of a regular-hours bar include the intervals of the
  overnight session before it. With each value average we form the same average of the availability indicator,
  $\bar a^{(k)}_{t,i}$, the share of the $k$ intervals on which the series was observed.
\item \textbf{Notation.} In tables and figures we drop the overbar and the interval index.
  \ifrawseries
  We refer to each exogenous series by its column name in the data (Table~\ref{tab:rh_dm_dictionary}): for instance,
  $\ser{vix}{\mathrm{VIX}}^{(25)}$ is the transformed VIX averaged over the 25 intervals before the forecast
  interval; the implied volatility series are \sername{vix}, \sername{vix3m} and \sername{vvix}; the suffixes
  \sername{_ewstock} and \sername{_vwstock} mark the equal- and value-weighted constituent aggregates (e.g.,
  \sername{ofi_ewstock} and \sername{spread_vwstock}); and the FOMC series are \sername{fomc_day},
  \sername{fomc_until_inv} and \sername{fomc_since_inv}. We write $\bar a^{(k)}(u)$ for the availability average of
  series $u$.
  \else
  For instance, $\mathrm{VIX}^{(25)}$ is the transformed VIX averaged over the 25 intervals before the forecast
  interval. We denote the implied volatility series by $\mathrm{VIX}$, $\mathrm{VIX3M}$ and $\mathrm{VVIX}$, label
  the equal- and value-weighted constituent aggregates with the superscripts EW and VW (e.g.,
  $\mathrm{OFI}^{\mathrm{EW}}$ and $\mathrm{Spread}^{\mathrm{VW}}$), denote the announcement predictors by
  $\mathrm{FOMC}$, and write $\bar a^{(k)}(u)$ for the availability average of series $u$.
  \fi
\item \textbf{HAR benchmark.} The HAR benchmark set holds the HAR components of the target, $\bar y^{(k)}_{t,i}$, the
  same trailing averages of the winsorized target $\tilde y$ (\rhdmNHorizonsPooled{} columns), and the calendar vector
  $\bm D_{t,i}$ (\rhorNCalendarAf{} columns): five day-of-the-week dummies and the hour of the interval-end stamp,
  entered as an integer, not as a set of dummies. The estimation sample holds no overnight interval, so
  $\bm D_{t,i}$ has no overnight indicator; and every interval in it falls on a weekday, so the five dummies sum to one
  and are not identified separately from the intercept: the OLS fit is the minimum-norm least-squares solution
  (Appendix~\ref{app:rh_methods_linear}).
\item \textbf{Full predictor set.} The full predictor set adds, for each of the \rhdmNSeriesPooled{} exogenous series
  of Section~\ref{sec:rh_data} (the FOMC group included), the standardized value averages (step~6 below) and the
  availability averages $\bar a^{(k)}$ at every horizon, so that it has
  \[
    \rhorNColsAf = \underbrace{\rhdmNHorizonsPooled}_{\bar y^{(k)}}
      + \underbrace{\rhorNCalendarAf}_{\bm D_{t,i}}
      + \underbrace{\rhdmNSeriesPooled \times \rhdmNHorizonsPooled}_{\bar u^{(k)}}
      + \underbrace{\rhdmNSeriesPooled \times \rhdmNHorizonsPooled}_{\bar a^{(k)}}
  \]
  candidate columns. The penalized linear models fit a screened subset of them, leaving out columns that are
  constant or exact duplicates on the estimation window, columns of a series with too few prints off its mode, and
  availability averages that move no more than missing observations at random would; each solver applies the
  screens on its own schedule (Appendix~\ref{app:rh_methods_linear}). The tree ensembles use every column.
\item \textbf{Transformations.} Each exogenous series passes through six steps. Every statistic in them is computed
  from earlier intervals only; Appendix~\ref{app:rh_dp_transform} gives the formulas and the constants.
  \begin{enumerate}
  \item \emph{Diurnal adjustment.} The raw series $u_{t,i}$ is divided by a time-of-day factor $b^{(u)}_{t,i}$: its
    mean (nonnegative series) or standard deviation (signed series) on interval $i$ over the prior days
    $\mathcal{D}_{t,i}$ of the target's factor $s^2_{t,i}$ (Section~\ref{sec:rh_methods_setup}), bounded relative to
    the factor's own trailing median. The implied volatility indices, the StockTwits sentiment score and the FOMC
    series are not adjusted.
  \item \emph{Variance stabilization.} A map chosen by what the series measures: the square root for variances,
    turnover and the effective spread, signed power roots for the higher moments and the autocovariance, the identity
    for the other signed series and absolute returns, and the logarithm otherwise. The volatility-demand series,
    whose standard-deviation divisor sits at its floor, take instead an inverse hyperbolic sine of the raw series,
    scaled by a trailing average of its absolute value.
  \item \emph{Winsorization.} A clip at rolling quantiles of the series' preceding \rhdmWinsorWindow{} intervals, at the
    levels of the target's winsorization (Section~\ref{sec:rh_methods_setup}). The bounds need a value on each of
    those intervals, so a series is not clipped where its prints leave a gap longer than the fill cap within the
    window: VIX3M and VVIX never are, VIX and the constituent aggregates on part of the sample. The FOMC series are
    not winsorized: the statement-day indicator passes through unchanged, and the proximity measures are
    deterministic ramps, each the extreme of its own trailing window between statements.
  \item \emph{Imputation.} An interval with $a_{t,i} = 0$ takes the median of the series' earlier observed
    transformed values, and before the series starts its first observed value; $a_{t,i}$ itself enters through
    $\bar a^{(k)}$.
  \item \emph{Trailing averages:} $\bar u^{(k)}_{t,i}$, $\bar a^{(k)}_{t,i}$ and $\bar y^{(k)}_{t,i}$ for
    $k \in \mathcal{K}$.
  \item \emph{Robust standardization.} Each value average is centred and scaled on the rolling window
    $\mathcal{W}_{t,i}$, the $W$ rows of the estimation sample before $(t,i)$ (Section~\ref{sec:rh_methods_setup}),
    imputed rows included:
    \begin{gather*}
      z^{(k)}_{t,i} = \frac{\bar u^{(k)}_{t,i}
          - \operatorname{med}\bigl\{\bar u^{(k)}_{t',i'} : (t',i') \in \mathcal{W}_{t,i}\bigr\}}
        {\max\bigl\{\mathrm{IQR}^{(k)}_{t,i},\ \rhdmIqrFloorFrac\,\overline{\mathrm{IQR}}^{(k)}_{t,i},\
          \mathrm{IQR}^{\ast(k)}_{t,i}\bigr\}}, \\
      \mathrm{IQR}^{(k)}_{t,i} = \mathrm{IQR}\bigl\{\bar u^{(k)}_{t',i'} : (t',i') \in \mathcal{W}_{t,i}\bigr\},
    \end{gather*}
    where $\overline{\mathrm{IQR}}^{(k)}_{t,i}$ is the largest $\mathrm{IQR}^{(k)}$ up to $(t,i)$ and
    $\mathrm{IQR}^{\ast(k)}_{t,i}$ a causal floor, the interquartile range of the average over the earlier rows whose
    $k$ intervals were all observed. The median and the interquartile range are recomputed at every row; the rows of
    the first window, all before the first forecast, take that window's statistics
    (Appendix~\ref{app:rh_dp_scale}).
    \begin{itemize}
    \item The availability averages are not standardized; the HAR components and the calendar columns follow the same
      rule with $\mathrm{IQR}^{\ast(k)} = 0$.
    \item \emph{Late start.} For a series whose first print comes after the start of the sample,
      $z^{(k)}_{t,i} = 0$ on every row whose window $\mathcal{W}_{t,i}$ holds a row whose $k$-interval average
      reaches back before the day of that print. Because the window spans \rhorWindowBars{} bars (about
      \rhorWindowDays{} trading days), late series enter late: on the bars scored from \rhorSpanStart{} to
      \rhorSpanEnd, the value averages are live on \rhlvVixEvRange\% of bars for the VIX, \rhlvVixThreeEvRange\% for
      VIX3M, \rhlvStAttEvRange\% for StockTwits attention, \rhlvStSentEvRange\% for its sentiment and tag count,
      \rhlvVdEvRange\% for volatility demand and \rhlvVvixEvRange\% for VVIX (ranges over $k$).
    \item \emph{No observation.} $z^{(k)}_{t,i} = 0$ whenever the $k$-interval average holds no observed interval,
      $\bar a^{(k)}_{t,i} = 0$, so the averages of a feed that has stopped are neutral.
    \item The step equalizes interquartile ranges, not variances, and scales each row on its own trailing window;
      the penalized linear models therefore standardize the columns again within their estimation window
      (Appendix~\ref{app:rh_methods_linear}).
    \end{itemize}
  \end{enumerate}
\item \textbf{No look-ahead.} Every statistic indexed by $(t,i)$ uses intervals before $(t,i)$ only: the trailing
  averages, the diurnal factors and their bounds, the winsorization bounds, the imputation medians, the routing of
  step~2 and the standardization statistics. The FOMC proximity measures use the pre-announced calendar
  (Section~\ref{sec:rh_data}). Two inputs use later intervals, and neither looks ahead of a forecast: the first print
  that fills the intervals before a series starts (step~4), which the late-start rule sets to zero, and the statistics
  that scale the rows of the first standardization window, all of which precede the first forecast.
\end{itemize}

